\section{Non-perturbative renormalization conditions}
\label{sec:idea}

In this section, the non-perturbative renormalization scheme
is presented in detail.  We also
discuss the conditions  which must be satisfied for our method to be
applicable.

To facilitate the discussion, it is convenient to classify composite operators
into three main classes, according to their ultra-violet behaviour, as
$a \to 0$:
\begin{enumerate}
\item {\bf Finite operators}: these include the vector and axial vector
currents. Another example  is  the ratio
of the  scalar and pseudoscalar renormalization constants, $Z_S/Z_P$. For
these cases, the Ward identity method is also applicable.
\item {\bf Logarithmically divergent operators}:
this is a large family of operators, relevant to hadron phenomenology.
It includes
   scalar and pseudoscalar densities,  four-fermion $\Delta F=2$
operators, some components of the energy-momentum tensor, etc.
\item
 {\bf Power divergent operators}: these operators are present
when  mixing with lower dimensional
operators is possible. This happens in regularizations which have  an intrinsic
mass scale, such as  the lattice regularization.
Examples are often encountered in phenomenological applications of
lattice $QCD$, such
as  four-fermion operators relevant to
$\Delta I=1/2$ transitions or operators of order $1/m$
in the heavy quark effective theory (HQET) \cite{effi}--\cite{efff}.
\end{enumerate}

For simplicity,
we first consider the two-fermion operators of class 2., and extend
the discussion to other operators in this class, and also to those
in classes
1. and 3., in the next section.
Thus, the formulae given below apply directly to the  pseudoscalar
and scalar densities. Throughout our discussion, we assume that discretization
errors are negligible: in field theory language, this is equivalent
to the statement that
the renormalized Green functions do not differ appreciably from their
values in the limit of an infinite  ultra-violet cut-off.
The discussion is presented in the  limit of small quark masses and
in  Euclidean space-time. The discretisation of the quark action is assumed
to be performed {\it \`a la} Wilson, characterized by   explicit chiral
symmetry
breaking of the Lagrangian\footnote{This includes the SW-Clover action
\cite{slami}.}. The extension to staggered fermions
 is straightforward.

Let us consider the forward amputated Green function  $\Gamma_{O}(pa)$,
of a two-fermion bare lattice  operator $O(a)$, computed between off-shell
quark
states  with four-momentum $p$, with $p^2=\mu^2$, and
in a fixed gauge, for example the Landau gauge.
Without gauge-fixing all Green functions
computed between quark and gluon external
states are zero. This is  also true when the gauge is not
completely fixed, e.g.\ in the Coulomb gauge.
We define the renormalized operator $O(\mu)$, by introducing
the renormalization constant $Z_O$
\beq O(\mu)= Z_O\Bigl(\mu a, g(a)\Bigr) O(a). \label{eq:zren} \eeq
$Z_O$ is found, by imposing the renormalization condition
\beq   Z_O\Bigl(\mu a,  g(a)\Bigr) Z^{-1}_\psi
\Bigl(\mu a,  g(a)\Bigr)
\Gamma_{O}(pa)\vert_{p^2=\mu^2}=1, \label{eq:zdef} \eeq
where $Z_\psi$ is the field renormalization constant, to be defined below.
This  procedure defines a renormalized operator $O(\mu)$
which is  independent of the regularization scheme
\cite{altarelli}--\cite{ciu}.
 It depends,  however, on the external states and on the gauge. This does not
affect the final results, which, combined with the continuum calculation
of the renormalization conditions, at any given order of perturbation
theory,
will be gauge invariant
and independent of the external states.
Let us specify the different quantities entering  eq.\ (\ref{eq:zdef}).
$\Gamma_{O}$ is defined in terms of the expectation value\footnote{
 ``Expectation value" means, as  usual, that one averages the Green functions
 over the gauge field configurations, generated by  Monte Carlo
 simulation, see sec.\ \ref{sec:duef}.}  of
the non-amputated Green function $G_{O}(pa)$, and
of the quark propagator $S(pa)$
\beq \Gamma_{O}(pa)= \frac{1}{12} {\rm tr} \left( \Lambda_{O}(pa)\hat  P_O
 \right), \label{eq:proj} \eeq
where
\beq \Lambda_{O}(pa)=  S(pa)^{-1} G_{O}(pa) S(pa)^{-1}. \label{eq:lam}
\eeq
$\hat P_O$ is a suitable projector on the tree-level operator: $\hat P_O=
\hat 1$ ($\hat P_O= \gamma_5$) for the scalar (pseudoscalar) density.
The factor $1/12$ ensures the correct overall normalization of the
trace (colour$\times$spin=12).
Projectors are very convenient when defining   Green functions,
particularly in the non-perturbative case. They
 have been extensively used in refs.\ \cite{buras,ciu}. Of course  one can
also use other definitions of $Z_O$.

$Z^{1/2}_\psi$ is the renormalization constant of the fermion field.
It can be defined in different ways, some of which are equivalent
perturbatively. Beyond perturbation theory, the most natural
definition of $Z_\psi$ is obtained from the amputated Green function
of the conserved vector current $V^{C}$. Indeed, one knows
that for $V^{C}$ the renormalization constant is equal to one:
\beq
Z^{-1}_{V^{C}}= \Gamma_{V^C} \times Z^{-1}_{\psi}=
\frac{1}{48} {\rm tr}
\left( \Lambda_{V^{C}_{\mu}} (pa)\gamma_{\mu}\right)\vert_{p^2=\mu^2} \times
Z^{-1}_{\psi}=1,
\eeq
which implies
\beq
Z_{\psi}=\frac{1}{48} {\rm tr}
\left( \Lambda_{V^{C}_{\mu}} (pa)\gamma_{\mu}\right)\vert_{p^2=\mu^2}.
\label{eq:zpsi}
\eeq

Equations (\ref{eq:zdef})--(\ref{eq:zpsi}) completely define our method.
In the remainder of this section, we discuss some important aspects
concerning its applicability.

In the above formulae, for simplicity,
we have always
considered only forward matrix elements. In general, one has the freedom
to define the renormalization conditions at different external
momenta $p$ and $p^\prime$ ($p \neq p^\prime$).
The virtualities of the  quark states must be  much larger than
$\Lambda_{QCD}$.
 The reason is that, in
order to obtain the physical result, we have to combine the matrix element
of the renormalized operator $O(\mu)$, with a Wilson coefficient function.
The latter is computed
in continuum perturbation theory, by expanding in
$\as$ at a scale of order $\mu$. Thus,
for the validity of this perturbative calculation, $\mu$ must be large.
An important question in our program is whether it is possible to find, on
the lattice, a scale $\mu$ which is sufficiently low, in order to
have small  $O(a)$ effects,
and sufficiently large, in order to have small higher order corrections.
The range of $\mu$, where the above conditions are satisfied, depends
on the value of $g_0^2$ of the numerical simulation.
We stress that, if such a  window
does not exist, in current lattice simulations,
an accurate matching of  lattice operators to
 continuum ones becomes impossible,
not only in our approach, but also with any other method.

It might be expected that the condition $\mu \gg \Lambda_{QCD}$ ensures
that perturbation
theory is valid. When
spontaneous symmetry breaking occurs however, as  is the case in $QCD$, a large
value of $\mu$ may  not be enough, because of the
presence of  the Goldstone boson, the pion in our case.
At low momentum transfers $q=
p-p^\prime$,  Green functions  can receive  a
non-perturbative contribution from
the pion pole, which cannot be computed. The na\"\i ve expectation is that this
 effect vanishes
 as $1/q^2$, at large $q^2$. However, with fermion fields, this
contribution is proportional to  $1/p^2=1/\mu^2$, even when $q^2=0$.
The proof is given in the appendix.
We conclude that a large value of $\mu$ solves the problem. At low values
of $\mu^2$, however, for the pseudoscalar density
and the axial current, we expect to find visible  effects,
see secs.\ \ref{sec:casivari} and \ref{sec:nonper}.

A  danger to the non-perturbative scheme may come from the
presence of  Gribov copies. We do not address  this problem here;
it  has been investigated in refs.\ \cite{altri,gribov}.
The results of
ref.\ \cite{gribov} indicate that,
at least for the quantities of interest in the present study,  the ``Gribov
uncertainty'', in current lattice simulations, is at most of the same order as
 the statistical error.

One may imagine applying the renormalization conditions at very
large values of $\beta$, on lattices which have small physical volumes.
This might seem to overcome the problem of the existence of the window in
$\mu$.
For  physical applications however, it is necessary to evolve the
renormalization constants to smaller values of $\beta$.
Non-perturbative contributions, which are not detected at large $\beta$'s,
may then become important. This demonstrates the necessity,
for any  method,
of the existence of a  window $\Lambda_{QCD} \ll \mu \ll 1/a$, where
$a$ is the lattice spacing used in  numerical simulations of physical
quantities.  Only then are  non-perturbative effects and  lattice artefacts
small simultaneously.
